\scanpage{114}
\mspage{56}
\margin{ceci marche au moins dans la Cat des espaces analytiques}
est un mono $\leftarrow$, et (167) est déterminé par la condition que son composé avec (168) est donné par les morphismes
\diagram{\begin{tikzcd}[ampersand replacement=\&, column sep=large,row sep=small]
\& D_d\\
D_{d,d',d''}\arrow[ur]\arrow[dr] \&\\
\& D_{d''}
\end{tikzcd}}
qui figurent dans (\unsure{166}).

\margin{on pourrait l'appeler aussi ``\uline{structure de drapeaux} analytiques''}
Une ``\uline{Structure de déploiement d'\unsure{intersection} normale}'' se définit \inserted{d'abord} comme une structure de ``catégorie \deleted{multiplicité} analytique'' $\mathcal D$, i.e. d'une catégorie dans la catégorie des espaces analytiques
\margin{(d'abord \unsure{$\mathcal E$} dans la cat des \uline{espaces} \unsure{analytiques})}
\diagram{
(169)\quad\left\{\begin{array}{ll}
D_* & \text{espace analytique des objets}\\[5pt]
D_{**} & \text{\rule{4cm}{0.3pt}\quad flèches}\\[8pt]
\begin{tikzcd}[ampersand replacement=\&,row sep=small,column sep=small,baseline=(current bounding box.center)]
\& D_{**}\arrow[dl,"\unsure{\sigma|s}"']\arrow[dr,"b"]\&\\
D_*\&\&D_*
\end{tikzcd}
&\text{morphismes ``source'' et ``but''}\\[8pt]
D_*\longrightarrow D_{**} & \text{morphisme ``\unsure{identique}''}\\[8pt]
\multicolumn{2}{l}{(D_{**},b)\times_{D_*}(D_{**},\unsure{\sigma|s})\longrightarrow D_{**}}\\[5pt]
\multicolumn{2}{c}{(u,v)\longmapsto v\circ u}\\[5pt]
\multicolumn{2}{l}{\text{(morphisme ``composition des flèches'')}}
\end{array}\right.
}
\deleted{Et} Ces données satisfont aux axiomes habituels (associativité, unité\ldots) exprimant que
\margin{113}
